\documentclass[11pt,a4paper]{article}

\usepackage[utf8]{inputenc}
\usepackage[T1]{fontenc}
\usepackage[spanish,es-noshorthands]{babel}
\usepackage{lmodern}
\usepackage[a4paper,margin=2.5cm]{geometry}
\usepackage{amsmath,amssymb,amsthm,mathtools}
\usepackage[dvipsnames]{xcolor}
\usepackage[most]{tcolorbox}
\usepackage{enumitem}
\usepackage{titlesec}
\usepackage[colorlinks=true,linkcolor=RoyalBlue,citecolor=RoyalBlue]{hyperref}

% --- Estilo -----------------------------------------------------------------
\numberwithin{equation}{section}
\setlist[itemize]{label={\color{gray!85}\raisebox{0.32ex}{\rule{3.4pt}{3.4pt}}}, leftmargin=1.8em, itemsep=3pt}
\titleformat{\section}{\large\bfseries}{\thesection.}{1em}{}
\titleformat{\subsection}{\normalsize\bfseries}{\thesubsection}{1em}{}
\setlength{\parindent}{0pt}
\setlength{\parskip}{6pt}

\newtcolorbox{hilo}{colback=gray!7,colframe=gray!30,boxrule=0.5pt,arc=1pt,
  left=8pt,right=8pt,top=5pt,bottom=5pt,fontupper=\small}

% --- Notación ---------------------------------------------------------------
\newcommand{\vc}[1]{\vec{#1}}
\newcommand{\nab}{\vec{\nabla}}
\newcommand{\pd}[2]{\frac{\partial #1}{\partial #2}}
\newcommand{\ux}{\hat u_x}
\newcommand{\uy}{\hat u_y}
\newcommand{\uz}{\hat u_z}
\newcommand{\uR}{\hat u_R}

\title{\vspace{-1.5cm}\textbf{Gradiente de $R^{n}$}\\[4pt]
  \large\normalfont Regla de la cadena, por qué $\nab=-\nab'$ y su uso en las
  cargas de polarización}
\author{}
\date{\vspace{-0.8cm}\small Apuntes del 20 de octubre de 2026}

\begin{document}
\maketitle

\begin{hilo}
\textbf{Hilo conductor.} Todo descansa en una observación: $R^{n}$ depende de
las seis coordenadas $(x,y,z,x',y',z')$ \emph{solo a través de la distancia}
$R$, y $R$ depende de ellas solo a través de las \emph{diferencias}
$x-x'$, $y-y'$, $z-z'$. De lo primero sale (i) que el gradiente es radial y se
calcula con la regla de la cadena; de lo segundo, (ii) que derivar respecto
del punto fuente solo cambia el signo. El caso $n=-1$ es el paso del potencial
de Coulomb $1/R$ a su campo $\vc{R}/R^{3}$; su versión con $\nab'$,
$\nab'(1/R)=\uR/R^{2}$, es la llave que convierte el potencial de un material
polarizado en el de unas cargas $\rho_P$ y $\sigma_P$
(sección~\ref{sec:polarizacion}).
\end{hilo}

% ===========================================================================
\section{Planteamiento}
% ===========================================================================
El punto fuente es $\vc{r}\,'=(x',y',z')$ y el punto campo $\vc{r}=(x,y,z)$.
El vector relativo, su módulo y su versor son
\begin{align}
  \vc{R}&=\vc{r}-\vc{r}\,'=R_x\,\ux+R_y\,\uy+R_z\,\uz,
  \qquad R_x=x-x',\; R_y=y-y',\; R_z=z-z', \label{eq:Rvec}\\
  R&=\bigl[R_x^{2}+R_y^{2}+R_z^{2}\bigr]^{1/2},\qquad
  \uR=\frac{\vc{R}}{R}\quad\bigl(\vc{R}=R\,\uR\bigr). \label{eq:Rmod}
\end{align}
Queremos demostrar que, para cualquier $n$ real y $R\neq0$,
\begin{equation}\label{eq:objetivo}
  \nab(R^{n})=-\nab'(R^{n})=n\,R^{n-1}\,\uR=n\,R^{n-2}\,\vc{R}.
\end{equation}

% ===========================================================================
\section{Las piezas}
% ===========================================================================
\subsection{$R^{n}$ es una función compuesta}
Definimos $h(u)=u^{n}$. Entonces $R^{n}=h\bigl(R(x,y,z)\bigr)$: la dependencia
en $x$ pasa \emph{entera} por $R$. La regla de la cadena da
\begin{equation}\label{eq:cadena}
  \pd{R^{n}}{x}=h'(R)\,\pd{R}{x}=\pd{R^{n}}{R}\,\pd{R}{x}.
\end{equation}
El problema se reduce a dos derivadas de una variable.

\subsection{Derivada exterior: $\partial R^{n}/\partial R$}
Es la regla de la potencia, $\frac{\mathrm d}{\mathrm du}u^{n}=nu^{n-1}$:
\begin{equation}\label{eq:exterior}
  \pd{R^{n}}{R}=n\,R^{n-1}.
\end{equation}
Vale para \textbf{cualquier} $n\in\mathbb{R}$ (negativo, fraccionario\dots)
porque $R>0$. De ahí la hipótesis $R\neq0$:
\begin{itemize}
  \item si $n<1$, $R^{n-1}$ no está definido en $R=0$;
  \item además $R$ no es derivable en $\vc{r}=\vc{r}\,'$: su gráfica tiene allí
        un vértice en forma de cono (como $|x|$ en una dimensión).
\end{itemize}

\subsection{Derivada interior: $\partial R/\partial x$}
Derivamos $R=\bigl[(x-x')^{2}+(y-y')^{2}+(z-z')^{2}\bigr]^{1/2}$ respecto de
$x$, con todo lo demás fijo. Es otra regla de la cadena (raíz de un polinomio):
\begin{equation}\label{eq:dR}
  \pd{R}{x}
  =\tfrac12\bigl[\cdots\bigr]^{-1/2}\cdot 2(x-x')\cdot\underbrace{\pd{(x-x')}{x}}_{=\,1}
  =\frac{x-x'}{R}=\frac{R_x}{R}.
\end{equation}
\textbf{Interpretación.} $R_x/R$ es el coseno director de $\vc{R}$ con el eje
$x$: la distancia crece al mover $x$ tanto más cuanto más alineado esté
$\vc{R}$ con ese eje.

% ===========================================================================
\section{Demostración}
% ===========================================================================
\subsection{Componente $x$}
Sustituyendo \eqref{eq:exterior} y \eqref{eq:dR} en \eqref{eq:cadena} y
juntando potencias ($R^{n-1}/R=R^{(n-1)-1}$):
\begin{equation}\label{eq:compx}
  \pd{R^{n}}{x}=n\,R^{n-1}\,\frac{R_x}{R}=n\,R^{n-2}\,R_x .
\end{equation}

\subsection{Componentes $y$, $z$: simetría}
$R$ es simétrica en las tres coordenadas. Repitiendo \eqref{eq:dR} con $y$ o
$z$ se obtiene $\partial R/\partial y=R_y/R$ y $\partial R/\partial z=R_z/R$,
luego
\begin{equation}\label{eq:compyz}
  \pd{R^{n}}{y}=n\,R^{n-2}\,R_y,\qquad \pd{R^{n}}{z}=n\,R^{n-2}\,R_z .
\end{equation}

\subsection{Agrupando en el gradiente}
Por definición, $\nab f=\pd{f}{x}\ux+\pd{f}{y}\uy+\pd{f}{z}\uz$. El factor
$nR^{n-2}$ es un \emph{escalar común} a las tres componentes y sale fuera; lo
que queda entre paréntesis es exactamente \eqref{eq:Rvec}:
\begin{equation}\label{eq:agrupar}
  \nab(R^{n})=nR^{n-2}\bigl(R_x\,\ux+R_y\,\uy+R_z\,\uz\bigr)
  =nR^{n-2}\,\vc{R}.
\end{equation}
Finalmente, con $\vc{R}=R\,\uR$ de \eqref{eq:Rmod}:
\begin{equation}\label{eq:final}
  nR^{n-2}\,\vc{R}=nR^{n-2}\,R\,\uR=n\,R^{n-1}\,\uR .
\end{equation}

\subsection{La igualdad con $-\nab'$}
$\nab'$ deriva respecto de las coordenadas \emph{de la fuente}. Como $R$
depende de $x'$ solo a través de $x-x'$, la derivada interior cambia de signo:
\begin{equation}\label{eq:dRprima}
  \pd{R}{x'}=\tfrac12\bigl[\cdots\bigr]^{-1/2}\cdot2(x-x')\cdot
  \underbrace{\pd{(x-x')}{x'}}_{=\,-1}=-\frac{R_x}{R}.
\end{equation}
La derivada exterior \eqref{eq:exterior} no cambia, de modo que cada
componente de $\nab'(R^{n})$ es la de $\nab(R^{n})$ con signo opuesto:
$\nab'(R^{n})=-\nab(R^{n})$. Esto completa \eqref{eq:objetivo}.
\hfill$\blacksquare$

% ===========================================================================
\section{Comprobaciones e interpretación}
% ===========================================================================
\begin{itemize}
  \item \textbf{$n=1$:} $\nab R=\uR$. La distancia crece a ritmo $1$ al
        alejarse radialmente de la fuente, y no crece al moverse en
        perpendicular (sobre una esfera centrada en $\vc{r}\,'$).
  \item \textbf{$n=2$:} $\nab(R^{2})=2\vc{R}$. Se comprueba a mano:
        $\partial_x\bigl[(x-x')^{2}+\cdots\bigr]=2(x-x')$.
  \item \textbf{$n=-1$ (Coulomb):}
        \begin{equation}\label{eq:coulomb}
          \nab\!\left(\frac1R\right)=-\frac{1}{R^{2}}\,\uR=-\frac{\vc{R}}{R^{3}} .
        \end{equation}
        Con $V=\dfrac{q}{4\pi\varepsilon_0 R}$ y $\vc{E}=-\nab V$ se obtiene
        $\vc{E}=\dfrac{q}{4\pi\varepsilon_0}\dfrac{\vc{R}}{R^{3}}$: el campo de
        una carga puntual.
\end{itemize}

\textbf{Idea geométrica.} $R^{n}$ tiene simetría esférica alrededor de
$\vc{r}\,'$: es constante sobre cada esfera $R=\text{cte}$. El gradiente es
perpendicular a las superficies de nivel, así que \emph{tiene} que ser radial,
$\propto\uR$; y su módulo es la derivada en la dirección radial,
$\mathrm d(R^{n})/\mathrm dR=nR^{n-1}$. La cuenta en cartesianas solo confirma
esto componente a componente:
\begin{equation}\label{eq:general}
  \nab f(R)=f'(R)\,\uR,\qquad \nab' f(R)=-f'(R)\,\uR
  \qquad\text{para cualquier } f \text{ derivable}.
\end{equation}
La demostración de arriba es exactamente esta con $f(R)=R^{n}$: el único
ingrediente fue $\partial R/\partial x=R_x/R$, que no depende de $f$.

% ===========================================================================
\section{Aplicación: cargas de polarización}\label{sec:polarizacion}
% ===========================================================================
Un dieléctrico polarizado ocupa una región $V'$ limitada por la superficie
cerrada $S'=\partial V'$, de normal exterior $\hat n$. Su estado se describe con
la \emph{polarización} $\vc{P}(\vc{r}\,')$: el momento dipolar por unidad de
volumen. Queremos ver por qué su potencial es el de unas cargas
\begin{equation}\label{eq:rhoP-sigmaP}
  \phi(\vc{r})=\frac{1}{4\pi\varepsilon_0}\int_{V'}\frac{\rho_P}{R}\,\mathrm dV'
  +\frac{1}{4\pi\varepsilon_0}\oint_{S'}\frac{\sigma_P}{R}\,\mathrm dS',
  \qquad \rho_P=-\nab'\cdot\vc{P},\quad \sigma_P=\vc{P}\cdot\hat n .
\end{equation}

\subsection{Paso 1: del dipolo a $\nab'(1/R)$}
Un dipolo $\vc{p}$ en $\vc{r}\,'$ crea en $\vc{r}$ el potencial
$\phi=\dfrac{1}{4\pi\varepsilon_0}\dfrac{\vc{p}\cdot\uR}{R^{2}}$. El trozo
$\mathrm dV'$ del material es un dipolo $\mathrm d\vc{p}=\vc{P}\,\mathrm dV'$, y
sumando (superposición)
\begin{equation}\label{eq:pot-dip}
  \phi(\vc{r})=\frac{1}{4\pi\varepsilon_0}\int_{V'}
  \frac{\vc{P}\cdot\uR}{R^{2}}\,\mathrm dV'.
\end{equation}
Ahora entra \eqref{eq:objetivo}. Con $n=-1$ y en su versión con $\nab'$:
\begin{equation}\label{eq:nabla-prima-1R}
  \nab'\!\left(\frac1R\right)=-\nab\!\left(\frac1R\right)
  =-\bigl(-1\bigr)R^{-2}\,\uR=+\frac{\uR}{R^{2}} .
\end{equation}
Por tanto el integrando es
$\dfrac{\vc{P}\cdot\uR}{R^{2}}=\vc{P}\cdot\nab'\!\left(\dfrac1R\right)$.

\textbf{Por qué $\nab'$ y no $\nab$.} Integramos en $\vc{r}\,'$, que es donde vive
$\vc{P}$. Para mover derivadas de un factor a otro dentro de la integral (paso
2) y luego usar Gauss (paso 3), las derivadas tienen que ser respecto de la
\emph{variable de integración}. Con $\nab$ no se podría: $\vc{P}$ no depende
de $\vc{r}$.

\textbf{Signo.} $\uR$ apunta de la fuente al punto campo. Desplazar la
fuente $\vc{r}\,'$ en la dirección $+\uR$ la acerca a $\vc{r}$: $R$ decrece y
$1/R$ crece. Por eso $\nab'(1/R)$ apunta según $+\uR$, al revés que
$\nab(1/R)$.

\subsection{Paso 2: regla del producto (integración por partes en 3D)}
Para un escalar $f$ y un vector $\vc{A}$,
\begin{equation}\label{eq:producto}
  \nab'\cdot(f\vc{A})=f\,\nab'\cdot\vc{A}+\vc{A}\cdot\nab' f .
\end{equation}
Basta mirar una componente: $\partial_{x'}(fA_x)=f\,\partial_{x'}A_x+A_x\,\partial_{x'}f$
(regla del producto de una variable) y sumar las tres. Con $f=1/R$ y
$\vc{A}=\vc{P}$, y despejando el término que aparece en \eqref{eq:pot-dip}:
\begin{equation}\label{eq:despeje}
  \frac{\vc{P}\cdot\uR}{R^{2}}=\vc{P}\cdot\nab'\!\left(\frac1R\right)
  =-\frac{\nab'\cdot\vc{P}}{R}+\nab'\cdot\!\left(\frac{\vc{P}}{R}\right).
\end{equation}
La idea es la de integrar por partes: pasamos la derivada de $1/R$ a
$\vc{P}$ a cambio de una \emph{derivada total}, que se integrará en el borde.
Aquí se necesita que $\vc{P}$ sea de clase $C^{1}$ en $V'$.

\subsection{Paso 3: teorema de la divergencia}
Integrando \eqref{eq:despeje} en $V'$,
\begin{equation}\label{eq:integrado}
  \phi(\vc{r})=\frac{1}{4\pi\varepsilon_0}\left[
  \int_{V'}\frac{-\nab'\cdot\vc{P}}{R}\,\mathrm dV'
  +\int_{V'}\nab'\cdot\!\left(\frac{\vc{P}}{R}\right)\mathrm dV'\right].
\end{equation}
El teorema de la divergencia, $\int_{V'}\nab'\cdot\vc{F}\,\mathrm dV'
=\oint_{S'}\vc{F}\cdot\hat n\,\mathrm dS'$, aplicado a $\vc{F}=\vc{P}/R$, da
\begin{equation}\label{eq:gauss}
  \int_{V'}\nab'\cdot\!\left(\frac{\vc{P}}{R}\right)\mathrm dV'
  =\oint_{S'}\frac{\vc{P}\cdot\hat n}{R}\,\mathrm dS' .
\end{equation}
Cada término tiene ya la forma de un potencial de Coulomb,
$\frac{1}{4\pi\varepsilon_0}\int\frac{(\text{carga})}{R}$: llamando $\rho_P$ y
$\sigma_P$ a los numeradores se obtiene \eqref{eq:rhoP-sigmaP}.

\subsection{Paso 4: el punto $R=0$ (si $\vc{r}\in V'$)}
Si el punto campo está dentro del material, $1/R$ explota en $\vc{r}\,'=\vc{r}$
y el teorema de la divergencia no se puede aplicar directamente. Se excluye una
bola $B_\varepsilon$ de radio $\varepsilon$ alrededor de $\vc{r}$, se repiten los
pasos 2--3 en $V'\setminus B_\varepsilon$ (allí todo es regular) y se hace
$\varepsilon\to0$. Llamando $M=\max|\vc{P}|$ y $M'=\max|\nab'\cdot\vc{P}|$:
\begin{itemize}
  \item \textbf{Integrales de volumen.} Los integrandos son $O(1/R^{2})$ y
        $O(1/R)$, y en esféricas $\mathrm dV'=R^{2}\,\mathrm dR\,\mathrm d\Omega$
        anula la singularidad. Por ejemplo,
        $\bigl|\int_{B_\varepsilon}\frac{\nab'\cdot\vc{P}}{R}\,\mathrm dV'\bigr|
        \le M'\!\int_0^{\varepsilon}\frac{4\pi R^{2}}{R}\,\mathrm dR=2\pi M'\varepsilon^{2}\to0$.
  \item \textbf{Superficie de la bola} (aparece como borde extra):
        $\bigl|\oint_{\partial B_\varepsilon}\frac{\vc{P}\cdot\hat n}{R}\,\mathrm dS'\bigr|
        \le M\,\frac{4\pi\varepsilon^{2}}{\varepsilon}=4\pi M\varepsilon\to0$.
\end{itemize}
Todo lo que se añade o quita al excluir la bola tiende a cero, y el
resultado se mantiene.

\subsection{Interpretación: cargas ligadas}
\begin{itemize}
  \item \textbf{$\sigma_P=\vc{P}\cdot\hat n$.} En el interior, la carga
        positiva de un dipolo se compensa con la negativa del vecino. En la
        superficie no hay vecino: «asoma» carga, positiva donde $\vc{P}$ sale
        ($\vc{P}\cdot\hat n>0$) y negativa donde entra.
  \item \textbf{$\rho_P=-\nab'\cdot\vc{P}$.} Si $\vc{P}$ no es uniforme, la
        compensación interior falla: donde las flechas de $\vc{P}$ divergen
        sale más carga positiva de la que entra y queda un defecto negativo.
        Si $\vc{P}$ es constante, $\rho_P=0$ y solo hay carga superficial.
  \item \textbf{Carga total nula.} Con $\mathrm d\vc{S}\,'=\hat n\,\mathrm dS'$
        y otra vez el teorema de la divergencia,
        \begin{equation}\label{eq:QP}
          Q_P=-\int_{V'}\nab'\cdot\vc{P}\,\mathrm dV'
          +\oint_{S'}\vc{P}\cdot\hat n\,\mathrm dS'
          =-\oint_{S'}\vc{P}\cdot\mathrm d\vc{S}\,'
          +\oint_{S'}\vc{P}\cdot\mathrm d\vc{S}\,'=0 .
        \end{equation}
        Es lo esperable: cada dipolo es neutro, y por tanto también cualquier
        conjunto de ellos.
\end{itemize}

\textbf{Resumen del enlace.} La única pieza «de cálculo» que hizo falta fue
\eqref{eq:nabla-prima-1R}, el caso $n=-1$ de \eqref{eq:objetivo} con $\nab'$.
Gracias a ella, el integrando dipolar se escribe como un gradiente respecto
de la fuente, y el resto es regla del producto más teorema de la divergencia.

\end{document}
